\chapter{}
\label{cap:3}

Este cap\'itulo documenta en detalle la primera etapa del pipeline de
an\'alisis introducido en el Cap\'itulo~\ref{cap:1}: la \textbf{gesti\'on de
instancias}. Se describe, de manera t\'ecnica y espec\'ifica, el conjunto de
fuentes externas de instancias de programaci\'on semidefinida (SDP)
configuradas en el proyecto, los scripts de Python responsables de
descargarlas y extraerlas, la estructura de directorios en la que se
almacenan los artefactos intermedios, y el modo en que la interfaz de
usuario orquesta la ejecuci\'on de dichos scripts como subprocesos.

\section{Estado del pipeline y alcance del cap\'itulo}
\label{sec:cap3-alcance}

El pipeline de gesti\'on de instancias fue reconstruido \'integramente
durante el desarrollo documentado en este cap\'itulo, reemplazando una
implementaci\'on anterior de \'unico script (\texttt{prepare.py} monol\'itico)
por cinco scripts independientes, cada uno con responsabilidad \'unica en
conformidad con el principio establecido en el Cap\'itulo~\ref{cap:2}:
\textbf{descarga} (\texttt{download.py}), \textbf{limpieza}
(\texttt{clean.py}), \textbf{extracci\'on} (\texttt{extract.py}),
\textbf{clasificaci\'on} (\texttt{sort.py}) y \textbf{preparaci\'on}
(\texttt{prepare.py}). Los cinco constituyen la implementaci\'on vigente y
verificada, con sus resultados contrastados num\'ericamente contra la
implementaci\'on de referencia original (Secci\'on~\ref{sec:cap3-converters}).

Junto a ellos se encuentra \texttt{list\_files.py}, el script de
descubrimiento de archivos en \'indices HTML, que la interfaz de usuario
invoca para resolver las fuentes de tipo \texttt{index}
(Secci\'on~\ref{sec:cap3-index}). Este script proviene de la implementaci\'on
anterior y se conserv\'o sin cambios, ya que su redise\~no no formaba parte
del alcance de este trabajo.

Un matiz importante de integraci\'on: al cierre de este cap\'itulo, la
interfaz de usuario (\texttt{renderer.ts}) orquesta \'unicamente
\texttt{clean.py}, \texttt{download.py} y \texttt{extract.py}
(Secci\'on~\ref{sec:cap3-ui}). Los dos scripts m\'as nuevos,
\texttt{sort.py} y \texttt{prepare.py}, est\'an completos, probados de forma
exhaustiva por l\'inea de comandos contra datos reales, y listos para
integrarse a esa misma secuencia, pero dicha integraci\'on qued\'o pendiente
como pr\'oximo paso.

\section{Fuentes de instancias: \texttt{config/instances/sources.json}}
\label{sec:cap3-sources}

El cat\'alogo de fuentes externas de instancias SDP se declara de forma
declarativa en el archivo de configuraci\'on
\texttt{app/config/instances/sources.json}, en conformidad con la
convenci\'on general de separaci\'on entre configuraci\'on y c\'odigo
establecida en el Cap\'itulo~\ref{cap:2}. Cada entrada del arreglo
\texttt{sources} describe una fuente mediante los campos \texttt{name},
\texttt{url}, \texttt{type}, \texttt{description}, \texttt{format} y, de
manera opcional, \texttt{verified}. El campo \texttt{type} distingue entre
dos naturalezas de URL, que determinan el comportamiento del script de
descarga descrito en la Secci\'on~\ref{sec:cap3-download}:

\begin{description}
    \item[\texttt{\textquotedbl archive\textquotedbl}.] La URL apunta
    directamente a un \'unico archivo comprimido descargable (por ejemplo,
    un \texttt{.tar.gz}).
    \item[\texttt{\textquotedbl index\textquotedbl}.] La URL apunta a una
    p\'agina HTML de listado de
    directorio, dentro de la cual es necesario descubrir program\'aticamente
    los enlaces a los archivos individuales antes de poder descargarlos.
\end{description}

La Tabla~\ref{tab:cap3-fuentes} resume las cuatro fuentes actualmente
configuradas.

\begin{table}[htbp]
\centering
\caption{Fuentes de instancias SDP configuradas en \texttt{sources.json}.}
\label{tab:cap3-fuentes}
\begin{tabular}{p{2.6cm}p{2cm}p{7.5cm}}
\toprule
\textbf{Fuente} & \textbf{Tipo} & \textbf{Contenido} \\
\midrule
SDPLIB 1.2 & \texttt{archive} & Biblioteca de referencia para SDP: control,
dise\~no estructural, relajaciones combinatorias. Formato \texttt{dat-s}. \\
Mittelmann Benchmarks & \texttt{index} & Instancias de escala media y
grande mantenidas por Hans Mittelmann; requiere descubrimiento de enlaces. \\
DIMACS Library & \texttt{archive} & Programas mixtos
semidefinidos-cu\'adraticos-lineales del s\'eptimo desaf\'io DIMACS,
convertidos de SeDuMi a formato SDPA. \\
Biq Mac Library & \texttt{archive} & Instancias de programaci\'on
cuadr\'atica binaria no restringida (BQP), convertidas a relajaciones SDP de
tipo \emph{max-cut}. \\
\bottomrule
\end{tabular}
\end{table}

El campo \texttt{verified}, cuando presente y verdadero, se refleja en la
interfaz de usuario mediante una insignia \emph{``100\% working''} junto al
nombre de la fuente (v\'ease Secci\'on~\ref{sec:cap3-ui}), indicando que la
descarga de dicha fuente ha sido probada exitosamente de forma manual.

\section{Estructura de directorios de datos de instancias}
\label{sec:cap3-directorios}

En correspondencia con la subdivisi\'on de \texttt{data/} por etapa
conceptual descrita en el Cap\'itulo~\ref{cap:2}, los artefactos generados
por el pipeline de gesti\'on de instancias se organizan bajo
\texttt{app/data/instances/} en cinco subdirectorios con sem\'antica
estrictamente diferenciada, uno por cada etapa del pipeline:

\begin{description}
    \item[\texttt{raw/}.] Cach\'e fiel de las descargas: contiene
    exactamente los archivos obtenidos de las fuentes externas, sin
    modificaci\'on alguna —incluyendo, cuando corresponde, el archivo
    comprimido original (\texttt{.tar.gz}, \texttt{.zip}, entre otros)—. Cada
    fuente ocupa su propio subdirectorio, nombrado seg\'un el nombre de la
    fuente saneado por la funci\'on \texttt{sourceSubdir} de
    \texttt{renderer.ts} (Secci\'on~\ref{sec:cap3-ui}). Este directorio nunca
    es modificado por \texttt{extract.py}, lo que permite volver a ejecutar
    la extracci\'on cuantas veces sea necesario sin requerir una nueva
    descarga.
    \item[\texttt{extracted/}.] Resultado de descomprimir y aplanar el
    contenido de \texttt{raw/}, organizado en subdirectorios hom\'onimos por
    fuente. Antes de que \texttt{sort.py} se ejecute, puede contener archivos
    que no son instancias (generadores, \texttt{README}, entre otros); despu\'es
    de \texttt{sort.py} (Secci\'on~\ref{sec:cap3-sort}), contiene
    exclusivamente archivos de instancia, sin subcarpetas anidadas ni
    archivos comprimidos.
    \item[\texttt{generators/}.] Scripts generadores de instancias
    (t\'ipicamente \texttt{.m} o \texttt{.py}) que algunas fuentes distribuyen
    junto a sus datos —por ejemplo, la biblioteca DIMACS incluye funciones
    MATLAB como \texttt{genbisect.m} para reproducir ciertas instancias
    param\'etricamente—. \texttt{sort.py} los extrae de \texttt{extracted/} y
    los deposita aqu\'i, en un subdirectorio hom\'onimo por fuente.
    \item[\texttt{extras/}.] Todo archivo distribuido junto a las instancias
    que no es ni una instancia ni un generador —documentaci\'on
    (\texttt{README}, \texttt{FORMAT}, \texttt{CHANGELOG}), licencias, u otros
    archivos sin extensi\'on reconocida—. Igual que \texttt{generators/}, es
    poblado por \texttt{sort.py} y organizado por subdirectorio de fuente.
    \item[\texttt{prepared/}.] Resultado final del pipeline: instancias
    normalizadas al formato SDPA sparse (\texttt{.dat-s}), producidas por
    \texttt{prepare.py} (Secci\'on~\ref{sec:cap3-prepare}) a partir de
    \texttt{extracted/}. Es el \'unico directorio de datos referenciado
    directamente por el resto de la aplicaci\'on, a trav\'es del cat\'alogo
    \texttt{config/instances/catalog.json}.
\end{description}

Los cinco directorios son vaciados de forma conjunta por
\texttt{clean.py} (Secci\'on~\ref{sec:cap3-clean}), lo que garantiza que
nunca coexistan artefactos de una etapa posterior correspondientes a una
descarga previa ya descartada de \texttt{raw/}.

\section{El script \texttt{download.py}}
\label{sec:cap3-download}

El script \texttt{app/scripts/instances/download.py} tiene una
responsabilidad \'unica y deliberadamente acotada: descargar un \'unico
archivo desde una URL hacia \texttt{data/instances/raw/}, sin realizar
ning\'un tipo de descompresi\'on ni interpretaci\'on de su contenido. Esta
separaci\'on estricta entre descarga y extracci\'on —impuesta como decisi\'on
de dise\~no expl\'icita durante el redise\~no del pipeline— permite que
ambas responsabilidades puedan ser probadas, depuradas y reejecutadas de
forma independiente.

\subsection{Interfaz de l\'inea de comandos}

En conformidad con las convenciones de Python descritas en el
Cap\'itulo~\ref{cap:2}, el script se parametriza \'integramente mediante
\texttt{argparse}:

\begin{lstlisting}[language=bash]
python download.py <url> [--subdir SUBDIR]
\end{lstlisting}

El argumento posicional \texttt{url} es obligatorio y corresponde a la
direcci\'on del archivo a descargar. El argumento opcional \texttt{--subdir}
especifica el subdirectorio de \texttt{raw/} en el que se guardar\'a el
archivo; cuando se omite, el archivo se guarda directamente en la ra\'iz de
\texttt{raw/}. El nombre de archivo local se deriva de la propia URL
mediante la funci\'on \texttt{\_filename\_from\_url}, que toma el \'ultimo
segmento de la ruta descartando par\'ametros de \emph{query string}.

\subsection{Protocolo de progreso}

Durante la descarga, efectuada mediante \texttt{urllib.request.urlretrieve},
el script reporta el avance por \texttt{stdout} en forma de l\'ineas con el
prefijo literal \texttt{PROGRESS}, ya sea como porcentaje
(\texttt{PROGRESS 42.7}) cuando el servidor informa el tama\~no total del
recurso mediante la cabecera \texttt{Content-Length}, o como kilobytes
descargados (\texttt{PROGRESS 1024KB}) en caso contrario. Este protocolo es
consistente con el patr\'on general de comunicaci\'on de eventos
\texttt{python-stdout} descrito en el Cap\'itulo~\ref{cap:1}: la funci\'on
\texttt{appendDownloadLog} de \texttt{renderer.ts} filtra expl\'icitamente
estas l\'ineas al renderizar la bit\'acora visible, evitando saturar el
registro con actualizaciones de porcentaje.

\section{Descubrimiento de archivos en fuentes de tipo \texttt{index}}
\label{sec:cap3-index}

Cuando la fuente seleccionada declara \texttt{\textquotedbl type\textquotedbl: \textquotedbl index\textquotedbl}, la URL no
identifica un archivo descargable sino una p\'agina de listado de
directorio. Previo a invocar \texttt{download.py}, la interfaz de usuario
invoca el script auxiliar
\texttt{app/scripts/instances/list\_files.py}, que:

\begin{enumerate}
    \item Descarga el HTML de la p\'agina de \'indice mediante
    \texttt{urllib.request}, identific\'andose con una cabecera
    \texttt{User-Agent} para evitar bloqueos por parte del servidor.
    \item Analiza el HTML con una subclase de
    \texttt{html.parser.HTMLParser} que captura \'unicamente las etiquetas
    \texttt{<a href=\textquotedbl...\textquotedbl>} cuyo destino termine en una extensi\'on
    reconocida como instancia o archivo comprimido (\texttt{.dat-s},
    \texttt{.dat-s.gz}, \texttt{.tar.gz}, \texttt{.tgz}, \texttt{.zip},
    \texttt{.gz}), descartando enlaces de navegaci\'on relativos
    (\texttt{?...}, \texttt{/...}).
    \item Resuelve cada enlace relativo a una URL absoluta mediante
    \texttt{urllib.parse.urljoin}.
    \item Emite por \texttt{stdout} un \'unico objeto JSON con la forma
    \texttt{\{\textquotedbl files\textquotedbl: [\{\textquotedbl name\textquotedbl: ..., \textquotedbl url\textquotedbl: ...\}, ...]\}}.
\end{enumerate}

La funci\'on \texttt{downloadIndexSource} de \texttt{renderer.ts} consume
esta salida, filtra los archivos cuyo nombre termina en \texttt{.dat-s} o
\texttt{.dat-s.gz}, y descarga cada uno invocando \texttt{download.py} de
forma secuencial, actualizando la insignia de la fuente con el progreso
\texttt{(i/N)}.

\section{El script \texttt{extract.py}}
\label{sec:cap3-extract}

El script \texttt{app/scripts/instances/extract.py} constituye la segunda
etapa del pipeline: transforma el contenido crudo de \texttt{raw/} en
archivos de instancia planos dentro de \texttt{extracted/}. A diferencia de
\texttt{download.py}, que opera sobre una \'unica URL por invocaci\'on, este
script no recibe argumentos obligatorios y recorre la totalidad de
\texttt{raw/} en una \'unica ejecuci\'on, admitiendo opcionalmente
\texttt{--subdir} para restringir el alcance a una sola fuente.

\subsection{Algoritmo de extracci\'on por fuente}

Para cada subdirectorio de \texttt{raw/} (correspondiente a una fuente), la
funci\'on \texttt{extract\_source} ejecuta la siguiente secuencia sobre el
subdirectorio hom\'onimo de destino en \texttt{extracted/}:

\begin{enumerate}
    \item \textbf{Primera pasada.} Recorre recursivamente los archivos de
    \texttt{raw/<fuente>/}. Todo archivo cuyo nombre finalice en una de las
    extensiones reconocidas como archivo comprimido
    (\texttt{.zip}, \texttt{.tar.gz}, \texttt{.tgz}, \texttt{.tar.bz2},
    \texttt{.tar}) se extrae directamente hacia
    \texttt{extracted/<fuente>/} mediante los m\'odulos est\'andar
    \texttt{zipfile} o \texttt{tarfile} (este \'ultimo con el filtro de
    seguridad \texttt{filter=\textquotedbl data\textquotedbl}, que impide la escritura de rutas
    fuera del directorio de destino). Todo archivo que \emph{no} sea un
    archivo comprimido se copia tal cual mediante \texttt{shutil.copy2},
    omiti\'endose si ya existe un archivo hom\'onimo en el destino.
    \item \textbf{Extracci\'on recursiva de archivos anidados.} Algunas
    fuentes (como la biblioteca DIMACS) empaquetan archivos comprimidos
    dentro de otros archivos comprimidos. Tras la primera pasada, el
    script repite en bucle la b\'usqueda de archivos comprimidos dentro de
    \texttt{extracted/<fuente>/}, extrayendo y eliminando cada uno
    encontrado, hasta que no quede ninguno.
    \item \textbf{Descompresi\'on de miembros \texttt{.gz} individuales.}
    La funci\'on \texttt{\_gunzip\_nested\_members} recorre
    \texttt{extracted/<fuente>/} en busca de archivos \texttt{.gz}
    individuales —caso t\'ipico de la biblioteca DIMACS, cuyo archivo
    \texttt{.tar.gz} contiene instancias comprimidas individualmente— y los
    descomprime en el mismo lugar mediante el m\'odulo \texttt{gzip},
    eliminando el archivo \texttt{.gz} original tras la conversi\'on.
    \item \textbf{Aplanado del directorio.} La funci\'on
    \texttt{\_flatten\_dir} mueve todo archivo ubicado en una subcarpeta de
    \texttt{extracted/<fuente>/} hacia la ra\'iz de dicho directorio,
    resolviendo colisiones de nombre mediante la adici\'on de un sufijo
    num\'erico incremental (\texttt{\_1}, \texttt{\_2}, \ldots), y elimina
    a continuaci\'on las subcarpetas que hayan quedado vac\'ias. El
    resultado es que \texttt{extracted/<fuente>/} nunca contiene
    estructura de directorios anidada.
\end{enumerate}

Es importante notar que \texttt{raw/<fuente>/} no se modifica en ning\'un
momento de este proceso: los archivos comprimidos originales permanecen
intactos, y toda operaci\'on destructiva (eliminaci\'on de archivos
comprimidos ya extra\'idos, resoluci\'on de colisiones) ocurre
exclusivamente dentro de \texttt{extracted/}.

\section{El script \texttt{sort.py}}
\label{sec:cap3-sort}

Algunas fuentes distribuyen, junto a las instancias propiamente dichas,
archivos adicionales de naturaleza distinta: la biblioteca DIMACS incluye
funciones MATLAB generadoras (\texttt{genbisect.m}, \texttt{generate\_hamming.m},
\texttt{genfap.m}, \texttt{gentorus.m}), y SDPLIB incluye documentaci\'on de
texto plano (\texttt{README}, \texttt{FORMAT}, \texttt{CHANGELOG}). El script
\texttt{app/scripts/instances/sort.py} clasifica el contenido de
\texttt{extracted/} para que solo queden instancias reales en dicho
directorio, moviendo el resto a \texttt{generators/} o \texttt{extras/}
(Secci\'on~\ref{sec:cap3-directorios}).

\subsection{Criterio de clasificaci\'on}

La funci\'on \texttt{classify} punt\'ua cada archivo, en primer lugar, por
su extensi\'on:

\begin{description}
    \item[Generador.] Extensi\'on \texttt{.m} o \texttt{.py} —c\'odigo fuente,
    nunca datos de instancia.
    \item[Instancia.] Extensi\'on \texttt{.dat-s}, \texttt{.dat},
    \texttt{.mat}, \texttt{.sdpa}, \texttt{.sdp} o \texttt{.sparse}
    —permanece en \texttt{extracted/} sin moverse.
    \item[Extra.] Cualquier otra extensi\'on, incluyendo la ausencia de
    extensi\'on. Este caso por omisi\'on es deliberado: captura
    autom\'aticamente \texttt{README}, \texttt{FORMAT}, \texttt{CHANGELOG} y
    cualquier archivo de documentaci\'on futuro sin necesidad de enumerar sus
    nombres expl\'icitamente.
\end{description}

Esta \'ultima regla por omisi\'on tiene una excepci\'on, necesaria porque
los grafos de max-cut de Biq Mac Library no tienen una extensi\'on
reconocible (Secci\'on~\ref{sec:cap3-biqmac}): para archivos de esa fuente
espec\'ifica, \texttt{classify} recurre a \texttt{looks\_like\_triplet}
—la misma funci\'on de detecci\'on por contenido reutilizada desde
\texttt{instance\_grouping.py} (Secci\'on~\ref{sec:cap3-grouping})— antes
de clasificar el archivo como extra, evitando que dichos grafos terminen
err\'oneamente en \texttt{extras/}.

Para cada subdirectorio de fuente en \texttt{extracted/}, los archivos
clasificados como generador o extra se mueven —no se copian— hacia
\texttt{generators/<fuente>/} o \texttt{extras/<fuente>/} respectivamente,
cre\'andose dichos subdirectorios de destino bajo demanda. El script es
idempotente: si se ejecuta nuevamente sobre un \texttt{extracted/} ya
clasificado, no encuentra archivos que reclasificar y no produce cambios.
Al igual que \texttt{extract.py}, no recibe argumentos obligatorios y
recorre la totalidad de \texttt{extracted/} en una \'unica ejecuci\'on,
admitiendo opcionalmente \texttt{--subdir} para restringir el alcance a una
sola fuente.

\section{El script \texttt{clean.py}}
\label{sec:cap3-clean}

El script \texttt{app/scripts/instances/clean.py} restablece el pipeline de
instancias a su estado inicial. Vac\'ia recursivamente el contenido de los
cinco subdirectorios de datos —\texttt{raw/}, \texttt{extracted/},
\texttt{generators/}, \texttt{extras/} y \texttt{prepared/}—, eliminando
primero todos los archivos y, a
continuaci\'on, todas las subcarpetas que hayan quedado vac\'ias, sin
eliminar los propios directorios ra\'iz (que se recrean mediante
\texttt{mkdir(parents=True, exist\_ok=True)} si no existieran). El script
no recibe argumentos y reporta por \texttt{stdout} la cantidad total de
archivos eliminados, seguido de una \'unica l\'inea \texttt{PROGRESS 100}
para se\~nalizar finalizaci\'on inmediata a la interfaz de usuario.

Este script se invoca opcionalmente al inicio del flujo de descarga
descrito en la Secci\'on~\ref{sec:cap3-ui}, cuando el usuario activa el
interruptor \emph{``limpiar antes de descargar''}, evitando as\'i que
instancias de una fuente descartada en una sesi\'on anterior persistan
mezcladas con las de la sesi\'on actual.

\section{Conversi\'on de instancias a SDPA: \texttt{format\_converters/}}
\label{sec:cap3-converters}

Las instancias en \texttt{extracted/} pueden llegar en distintos formatos
seg\'un su fuente: SDPLIB ya distribuye \texttt{.dat-s} (SDPA sparse); la
biblioteca DIMACS distribuye muchas de sus instancias como archivos
\texttt{.mat} en \textbf{formato SeDuMi} (las variables \texttt{A}, \texttt{b},
\texttt{c}, \texttt{K} de la representaci\'on est\'andar de MATLAB para
programas c\'onicos); y Biq Mac Library distribuye la matriz de costos de
sus problemas combinatorios en dos formatos propios de texto plano
(Secci\'on~\ref{sec:cap3-biqmac}). El subdirectorio
\texttt{app/scripts/instances/format\_converters/} contiene el c\'odigo que
traduce cada uno de estos formatos a SDPA sparse.

\subsection{Conversi\'on directa SeDuMi $\to$ SDPA}
\label{sec:cap3-sedumi-converter}

Para el formato SeDuMi, la conversi\'on se implementa en dos capas: una
implementaci\'on de referencia heredada, de terceros, y una implementaci\'on
propia que la sustituye.

\subsubsection{La implementaci\'on de referencia: \texttt{extern/}}

El subdirectorio \texttt{format\_converters/extern/} contiene, sin
modificaci\'on, dos archivos de Brian Borchers (autor de SDPLIB y de
SDPpack) que definen la transformaci\'on matem\'atica de referencia:

\begin{description}
    \item[\texttt{writesdp.m}.] Funci\'on MATLAB que traduce un problema en
    formato SeDuMi al formato de texto SDPpack (SQLP).
    \item[\texttt{sqlptosdpa.c}.] Programa en C que traduce SDPpack a SDPA
    sparse, y que —cr\'iticamente— es responsable de una transformaci\'on
    matem\'atica no trivial: SDPA sparse no tiene un tipo de bloque nativo
    para conos cuadr\'aticos (\emph{second-order cones}, SOC), de modo que
    cada cono se re-expresa como un bloque semidefinido adicional en forma
    de \textbf{matriz flecha} (\emph{arrow matrix}): para un cono de tama\~no
    $n$, la condici\'on $x_1 \geq \|x_{2:n}\|_2$ es equivalente a que la
    matriz $n \times n$ con $x_1$ en la diagonal, $x_j/2$ en la posici\'on
    $(1,j)$ y ceros en el resto sea semidefinida positiva, agregando
    restricciones de igualdad adicionales (con lado derecho cero) que fijan
    la diagonal de dicha matriz y anulan sus entradas interiores.
\end{description}

Encadenados, \texttt{writesdp.m} y \texttt{sqlptosdpa.c} bastar\'ian para
convertir SeDuMi a SDPA, pero requieren un archivo de texto SQLP
intermedio por instancia y, para ejecutar \texttt{sqlptosdpa.c}, un
compilador de C —una herramienta no requerida por ning\'un otro script del
proyecto—. Por esa raz\'on se implement\'o \texttt{sedumi\_to\_dats.m}: la
misma transformaci\'on matem\'atica, fusionada en un solo paso.

\subsubsection{\texttt{sedumi\_to\_dats.m}: conversi\'on directa}

\texttt{sedumi\_to\_dats.m} y sus cuatro funciones auxiliares
(\texttt{sedumiBlockSizes.m}, \texttt{sedumiMatrixEntries.m},
\texttt{sedumiConeEntries.m}, \texttt{writeSdpaSparse.m}) reimplementan
desde cero el mismo algoritmo de \texttt{extern/}, evitando por completo el
archivo SQLP intermedio y la dependencia de un compilador de C. Se eligi\'o
MATLAB —y no Python ni C— porque \texttt{.mat} es su formato nativo (lectura
sin bibliotecas externas) y porque sus operaciones dispersas
(\texttt{find}, \texttt{accumarray}) ya est\'an implementadas en c\'odigo
compilado equivalente en velocidad al de C, sin agregar dependencias nuevas
al proyecto. En conformidad con el l\'imite de $100$ l\'ineas por archivo
MATLAB (Cap\'itulo~\ref{cap:2}), la l\'ogica se reparte en cinco archivos de
responsabilidad \'unica: normalizaci\'on de tama\~nos de bloque
(\texttt{K.l}/\texttt{K.q}/\texttt{K.s}), extracci\'on de entradas de los
bloques semidefinidos, construcci\'on de la representaci\'on de matriz
flecha para los conos, y escritura del archivo \texttt{.dat-s} final.

Dos decisiones de dise\~no de \texttt{sedumi\_to\_dats.m} merecen
explicaci\'on expl\'icita porque no son evidentes a partir del c\'odigo de
\texttt{extern/}:

\begin{description}
    \item[Ausencia de cambio de signo.] \texttt{writesdp.m} invierte el
    signo del vector de costos (\texttt{c = -c}) porque SDPpack maximiza y
    SeDuMi minimiza; \texttt{sqlptosdpa.c}, al escribir el archivo SDPA,
    invierte el signo \emph{nuevamente}. Ambas inversiones se cancelan, de
    modo que el resultado final coincide con el \texttt{c} original de
    SeDuMi sin cambio de signo alguno —hecho que \texttt{sedumi\_to\_dats.m}
    aprovecha para omitir la inversi\'on por completo, ya que SDPA y SeDuMi
    comparten la misma convenci\'on ($\min c^\top x \text{ s.a. } Ax=b$).
    \item[Extracci\'on de bloques sparse-nativa.] \texttt{writesdp.m}
    reconstruye cada bloque semidef\'inido como una matriz densa
    ($\mathtt{reshape}$) por cada restricci\'on, uno a la vez. Para un bloque
    grande y disperso —el archivo \texttt{bm1.mat} de DIMACS tiene un
    \'unico bloque de $882\times882$ con $883$ restricciones— materializar
    esa reconstrucci\'on densamente exigir\'ia varios gigabytes de memoria.
    \texttt{sedumiMatrixEntries.m} evita esto por completo: usa
    \texttt{find} directamente sobre la submatriz dispersa para obtener solo
    las entradas no nulas, agrup\'andolas por posici\'on sim\'etrica
    ($(i,j)$/$(j,i)$) mediante \texttt{accumarray} en lugar de reconstruir
    la matriz completa.
\end{description}

\subsubsection{L\'imite de tama\~no para conos individuales}

La representaci\'on de matriz flecha de un cono de tama\~no $n$ requiere
$(n-2)(n-1)/2$ restricciones adicionales para anular sus entradas
interiores —cuadr\'atico en $n$, id\'entico en \texttt{extern/} y en
\texttt{sedumi\_to\_dats.m}—. Para la mayor\'ia de las instancias esto es
manejable, pero algunas instancias DIMACS de la familia \texttt{sched}
tienen un \'unico cono de tama\~no superior a $17\,000$, lo que requerir\'ia
generar del orden de $1.6\times10^8$ restricciones —un archivo \texttt{.dat-s}
de varios gigabytes, impr\'actico para los solvers del proyecto—.
\texttt{sedumiConeEntries.m} calcula este conteo \emph{antes} de intentar
construir dichas restricciones y, si excede un l\'imite de seguridad
($5\times10^6$), aborta la conversi\'on de esa instancia con un mensaje
explicando la causa, en lugar de agotar la memoria del sistema intent\'andolo.
Este l\'imite descarta dos instancias del cat\'alogo DIMACS
(\texttt{sched\_200\_100\_orig}, \texttt{sched\_200\_100\_scaled}) y reduce a
un caso manejable otras cuatro (\texttt{sched\_100\_100\_*},
\texttt{sched\_100\_50\_*}), seg\'un se reporta en la
Secci\'on~\ref{sec:cap3-prepare}.

\subsubsection{Validaci\'on num\'erica}

La correcci\'on de \texttt{sedumi\_to\_dats.m} se verific\'o comparando su
salida, entrada por entrada, contra la de \texttt{extern/writesdp.m} +
\texttt{sqlptosdpa.c} compilado, para instancias representativas de cada
combinaci\'on de tipos de bloque —solo SDP (\texttt{bm1.mat}, bloque \'unico
de $882\times882$), muchos conos peque\~nos (\texttt{nb.mat}, $793$ conos de
tama\~no $3$) y la combinaci\'on completa L+Q+S (\texttt{filtinf1.mat})—,
con coincidencia exacta salvo por la numeraci\'on de las filas de
restricciones extra (irrelevante para SDPA, que no exige un orden
particular). El mismo procedimiento se emple\'o para medir el desempe\~no:
sobre \texttt{nql60old.mat} ($3\,600$ conos de tama\~no $3$),
\texttt{sedumi\_to\_dats.m} tarda $0.46\,\mathrm{s}$ frente a
$2.39\,\mathrm{s}$ del camino de dos etapas —aproximadamente $5\times$ m\'as
r\'apido, principalmente porque \texttt{sqlptosdpa.c} realiza dos pasadas de
lectura car\'acter por car\'acter sobre el archivo de texto intermedio.

\subsection{\texttt{biqmac\_to\_dats.py}: BQP y max-cut}
\label{sec:cap3-biqmac}

Biq Mac Library no distribuye relajaciones SDP ya construidas como SeDuMi:
distribuye la matriz de costos $Q$ del problema combinatorio original, en
texto plano, bajo dos formas seg\'un la familia de instancias:

\begin{description}
    \item[Programaci\'on cuadr\'atica binaria sin restricciones (UBQP).]
    Archivos \texttt{.sparse} (y, para un subconjunto, un \texttt{.mat}
    redundante que —pese a la extensi\'on— es tambi\'en texto plano, no un
    binario real de MATLAB) con la matriz sim\'etrica $Q$ del problema
    $\max x^\top Q x$, $x\in\{0,1\}^n$, en formato disperso triangular
    superior \emph{incluyendo} la diagonal (el t\'ermino lineal).
    \item[Grafos de max-cut ponderados.] Archivos sin extensi\'on confiable
    —Biq Mac codifica el n\'umero de instancia como un sufijo num\'erico que
    \texttt{pathlib} tambi\'en interpreta como extensi\'on, de modo que
    \texttt{g05\_100.0} y \texttt{g05\_100.1} son dos instancias
    \emph{distintas}, no dos formatos de la misma— con la lista de aristas
    del grafo a cortar, en el mismo formato triangular superior pero
    \emph{sin} diagonal.
\end{description}

Como ambos formatos comparten exactamente la misma forma de archivo
(cabecera \texttt{n cuenta}, luego l\'ineas \texttt{i j valor}),
\texttt{biqmac\_to\_dats.py} los distingue por contenido, no por extensi\'on:
si alguna entrada tiene $i=j$, es UBQP; si ninguna, es un grafo de max-cut.
Ambos casos exigen una transformaci\'on matem\'atica distinta hacia SDPA,
ninguna de las cuales tiene contraparte en \texttt{extern/}:

\begin{description}
    \item[UBQP.] Se homogeneiza introduciendo una variable auxiliar $y_0=1$
    y sustituyendo $x_i=(1+y_i)/2$, de modo que $x^\top Q x = \tfrac{1}{4}
    y^\top \tilde{Q} y$ con $y\in\{-1,1\}^{n+1}$ y
    $\tilde{Q}=\begin{psmallmatrix}S & r^\top\\ r & Q\end{psmallmatrix}$,
    donde $r$ es el vector de sumas de fila de $Q$ y $S$ la suma total de
    sus entradas. El bloque SDP resultante tiene tama\~no $n+1$, con
    $C=-\tfrac{1}{4}\tilde{Q}$ (el signo recupera la maximizaci\'on dentro
    de la convenci\'on de minimizaci\'on de SDPA) y $n+1$ restricciones
    $\mathrm{diag}(Y)=e$.
    \item[Max-cut.] No requiere variable auxiliar: el bloque SDP tiene
    tama\~no $n$ directamente, con $C_{ij}=\tfrac{1}{4}w_{ij}$ para cada
    arista $(i,j)$ y $n$ restricciones $\mathrm{diag}(X)=e$. Esta f\'ormula
    se verific\'o contra \texttt{maxG11.dat-s} de SDPLIB —una instancia de
    max-cut ya conocida como correcta, con la misma estructura de bloque
    \'unico y coeficientes $\pm0.25$— antes de implementarla.
\end{description}

\subsubsection{Validaci\'on por consistencia num\'erica}

A diferencia de \texttt{sedumi\_to\_dats.m}, para este formato no existe una
implementaci\'on de referencia en el proyecto contra la cual comparar
salida byte a byte. La validaci\'on se hizo, en cambio, verificando que el
valor objetivo se preserve: para varias asignaciones binarias aleatorias
$x$ (UBQP) o de corte $y\in\{-1,1\}^n$ (max-cut), se calcul\'o el valor
combinatorio directamente a partir de $Q$ y se compar\'o contra $C\bullet
Y$ evaluado sobre la matriz $C$ realmente escrita en el \texttt{.dat-s}
generado, confirmando coincidencia exacta (salvo error de punto flotante)
en ambos formatos y en instancias de tama\~no y peso variables.

\section{El script \texttt{prepare.py}}
\label{sec:cap3-prepare}

El script \texttt{app/scripts/instances/prepare.py} constituye la \'ultima
etapa del pipeline: normaliza el contenido de \texttt{extracted/} al
formato SDPA sparse dentro de \texttt{prepared/}, y regenera el cat\'alogo
\texttt{config/instances/catalog.json} que el resto de la aplicaci\'on usa
para enumerar las instancias disponibles (Cap\'itulo~\ref{cap:2}).

\subsection{Agrupaci\'on por instancia y prioridad de formato}
\label{sec:cap3-grouping}

Una misma instancia puede existir en \texttt{extracted/} en m\'as de un
formato simult\'aneamente: la biblioteca DIMACS distribuye, para varias
instancias, tanto un archivo \texttt{.dat} en formato de grafo (el problema
combinatorio original) como un archivo \texttt{.mat} en formato SeDuMi (su
relajaci\'on SDP ya construida) bajo el mismo nombre base —por ejemplo,
\texttt{bm1.dat} y \texttt{bm1.mat}—; Biq Mac Library distribuye, para su
familia UBQP, tanto un \texttt{.sparse} real como un \texttt{.mat}
falsificado (texto plano, no un binario de MATLAB —
Secci\'on~\ref{sec:cap3-biqmac}) bajo el mismo nombre base. Esta l\'ogica de
agrupamiento y prioridad se factoriz\'o en su propio m\'odulo,
\texttt{app/scripts/instances/instance\_grouping.py}, en lugar de vivir
dentro de \texttt{prepare.py}, porque tambi\'en la reutiliza
\texttt{sort.py} (Secci\'on~\ref{sec:cap3-sort}) para reconocer los grafos
de Biq Mac que no tienen extensi\'on confiable.

La funci\'on \texttt{group\_by\_stem} agrupa por nombre base
(\texttt{Path.stem}) \'unicamente los archivos cuya extensi\'on es una de
las tres reconocidas (\texttt{.dat-s}, \texttt{.sparse}, \texttt{.mat}); el
resto se deja sin agrupar. Esta distinci\'on es necesaria porque agrupar
por \texttt{Path.stem} indiscriminadamente es incorrecto para los grafos de
Biq Mac: \texttt{pathlib} interpreta el sufijo num\'erico de instancia que
Biq Mac usa para distinguir corridas —\texttt{g05\_100.0},
\texttt{g05\_100.1}, \ldots— como si fuera una extensi\'on de archivo, de
modo que agruparlos por \texttt{stem} fusionar\'ia diez instancias
\emph{distintas} bajo un mismo grupo \texttt{g05\_100} en lugar de
reconocerlas como diez archivos independientes.

Con esa agrupaci\'on ya hecha, \texttt{plan\_source} decide, por grupo, una
\'unica acci\'on seg\'un la prioridad siguiente:

\begin{enumerate}
    \item Si el grupo incluye un \texttt{.dat-s}, se copia directamente
    (\texttt{shutil.copy2}) hacia \texttt{prepared/<fuente>/<nombre>.dat-s}
    sin invocar ning\'un conversor.
    \item En caso contrario, si el grupo incluye un \texttt{.sparse}, se
    encola para conversi\'on v\'ia \texttt{biqmac\_to\_dats.py}
    (Secci\'on~\ref{sec:cap3-biqmac}) —con prioridad expl\'icita sobre
    \texttt{.mat}, para que el \texttt{.mat} falsificado de Biq Mac nunca
    se env\'ie por error al conversor de SeDuMi.
    \item En caso contrario, si el grupo incluye un \texttt{.mat}, se
    encola para conversi\'on v\'ia \texttt{sedumi\_to\_dats.m}
    (Secci\'on~\ref{sec:cap3-sedumi-converter}).
    \item Para cada archivo que qued\'o sin agrupar (por extensi\'on no
    reconocida): si pertenece a Biq Mac Library y su contenido calza con
    el formato de tripletas \texttt{n cuenta} / \texttt{i j valor}
    (\texttt{looks\_like\_triplet}), se encola tambi\'en para
    \texttt{biqmac\_to\_dats.py}, usando el nombre de archivo completo
    —no el \texttt{stem}— como base del \texttt{.dat-s} de salida, para
    preservar la distinci\'on entre instancias como \texttt{g05\_100.0} y
    \texttt{g05\_100.1}.
    \item De lo contrario —como los archivos \texttt{.dat} de grafo DIMACS
    que carecen de contraparte \texttt{.mat} (\texttt{biomedP},
    \texttt{fap36}, \texttt{industry2})—, la instancia se omite y se
    reporta por \texttt{stdout} bajo el encabezado \emph{``Skipped \ldots\
    instance(s) with no supported format''}.
\end{enumerate}

\subsection{Una \'unica sesi\'on de MATLAB para todas las conversiones}

Cada arranque de MATLAB tiene un costo fijo de varios segundos; lanzar un
proceso de MATLAB por cada archivo \texttt{.mat} —como har\'ia invocar
\texttt{sedumi\_to\_dats.m} una vez por instancia desde Python, seg\'un el
patr\'on \texttt{subprocess} ya usado para \texttt{buildIS}
(Cap\'itulo~\ref{cap:1})— multiplicar\'ia ese costo por decenas de
instancias. En su lugar, \texttt{build\_batch\_script} genera un \'unico
script MATLAB temporal que antepone un \texttt{addpath} hacia
\texttt{format\_converters/} y encadena un bloque \texttt{try/catch} por
cada conversi\'on pendiente —de modo que el fallo de una instancia (por
ejemplo, por el l\'imite de tama\~no de cono de la
Secci\'on~\ref{sec:cap3-converters}) no interrumpe las dem\'as—, y
\texttt{run\_matlab\_jobs} lo ejecuta con \texttt{matlab -batch
run('<script>')} en un \'unico proceso, transmitiendo su salida por
\texttt{stdout} l\'inea por l\'inea igual que el resto de los scripts del
pipeline.

\subsection{Regeneraci\'on del cat\'alogo}

Al finalizar la copia y las conversiones, \texttt{write\_catalog} recorre
recursivamente \texttt{prepared/} en busca de archivos \texttt{.dat-s} y
escribe \texttt{config/instances/catalog.json} con una entrada por
instancia (\texttt{name}, \texttt{file}, \texttt{source}, \texttt{path}),
reemplazando por completo el contenido anterior del archivo. Esta
regeneraci\'on total —en lugar de una actualizaci\'on incremental— garantiza
que el cat\'alogo jam\'as referencie una instancia que ya no existe en
\texttt{prepared/}, a costa de exigir que \texttt{prepare.py} se ejecute
sobre la totalidad del pipeline en cada corrida.

\section{Orquestaci\'on desde la interfaz de usuario}
\label{sec:cap3-ui}

La p\'agina \emph{Instances $\to$ Download} de la interfaz, implementada en
\texttt{app/src/renderer/renderer.ts}, no contiene l\'ogica de descarga ni
de extracci\'on propia: su responsabilidad se limita a invocar, mediante
\texttt{window.api.runPython} (Cap\'itulo~\ref{cap:1}), los scripts descritos
en este cap\'itulo en la secuencia correcta, y a reflejar su progreso en la
interfaz gr\'afica a trav\'es de tarjetas de fuente
(\texttt{source-card}) con insignia de estado (\emph{Idle},
\emph{Downloading}, \emph{Done}, \emph{Error}) y una bit\'acora de ejecuci\'on
en tiempo real (\texttt{download-log}).

La funci\'on \texttt{downloadSources}, invocada al presionar el bot\'on
\emph{``Download selected''}, ejecuta la secuencia ilustrada en la
Figura~\ref{fig:cap3-flujo}:

\begin{enumerate}
    \item Si el interruptor de limpieza est\'a activo, ejecuta
    \texttt{clean.py} y espera su finalizaci\'on antes de continuar.
    \item Para cada fuente seleccionada por el usuario, en orden: si la
    fuente es de tipo \texttt{index}, delega en
    \texttt{downloadIndexSource} (Secci\'on~\ref{sec:cap3-index}); en caso
    contrario, invoca \texttt{download.py} una \'unica vez con la URL de la
    fuente. El usuario puede interrumpir el proceso en cualquier momento
    mediante el bot\'on \emph{``Stop''}, verificado entre cada fuente y entre
    cada archivo individual de una fuente de tipo \texttt{index}.
    \item Una vez procesadas todas las fuentes seleccionadas (o
    interrumpido el proceso), si no hubo interrupci\'on, ejecuta
    \texttt{extract.py} sin argumentos —una \'unica vez, sobre la totalidad
    de \texttt{raw/}— y actualiza el indicador de estado global
    (\texttt{pipeline-status}) a \emph{``Ready''} o \emph{``Extract
    failed''} seg\'un el c\'odigo de salida del proceso.
\end{enumerate}

\begin{figure}[htbp]
\centering
\begin{tikzpicture}[
    box/.style={draw, rounded corners, align=center, minimum height=10mm,
        minimum width=30mm, font=\small, inner sep=2mm},
    ui/.style={box, fill=blue!8},
    py/.style={box, fill=orange!12},
    fs/.style={box, fill=green!10},
    arr/.style={-{Stealth[length=2.2mm]}, thick},
    darr/.style={-{Stealth[length=2mm]}, thick, dashed},
    lbl/.style={font=\scriptsize, align=center, fill=white, inner sep=0.5pt},
    stepnum/.style={lbl, font=\scriptsize\bfseries}
]
    % fila superior: interfaz
    \node[ui] (ui) at (6,9.4) {renderer.ts\\\texttt{downloadSources()}};

    % fila de scripts (orden de ejecucion, izquierda a derecha)
    \node[py] (clean)    at (2,7.0) {\texttt{clean.py}};
    \node[py] (download) at (6,7.0) {\texttt{download.py}\\\footnotesize(por fuente)};
    \node[py] (extract)  at (10,7.0) {\texttt{extract.py}\\\footnotesize(al final)};

    % script auxiliar de descubrimiento (solo fuentes tipo index)
    \node[py] (listfiles) at (4,4.6) {\texttt{list\_files.py}\\\footnotesize(solo tipo \texttt{index})};

    % fila de directorios de datos
    \node[fs] (prepdir) at (2,2.2) {\texttt{data/instances/}\\\texttt{prepared/}};
    \node[fs] (rawdir)  at (6,2.2) {\texttt{data/instances/}\\\texttt{raw/}};
    \node[fs] (extdir)  at (10,2.2) {\texttt{data/instances/}\\\texttt{extracted/}};

    % secuencia de ejecucion: 1 clean, 2 download, 3 extract
    \draw[arr] (ui.west) -- (clean.north) node[stepnum, pos=0.65, left] {1};
    \draw[arr] (clean.east) -- (download.west) node[stepnum, midway, above] {2};
    \draw[arr] (download.east) -- (extract.west) node[stepnum, midway, above] {3};

    % list_files.py alimenta a download.py con URLs descubiertas
    \draw[arr] (listfiles.north) -- (download.south) node[lbl, midway] {JSON};

    % escritura de resultados
    \draw[arr] (download.south) -- (rawdir.north) node[lbl, midway] {escribe};
    \draw[arr] (extract.south) -- (extdir.north) node[lbl, midway] {escribe};

    % clean.py vacia los tres directorios de datos (bus inferior, bajo la fila de datos)
    \draw[darr] (clean.south) -- (prepdir.north) node[lbl, midway, left] {vac\'ia};
    \draw[darr] (clean.south) -- (2,0.5) -- (6,0.5) -- (rawdir.south);
    \draw[darr] (clean.south) -- (2,0.15) -- (10,0.15) -- (extdir.south);
\end{tikzpicture}
\caption{Secuencia de invocaci\'on de los scripts del pipeline de instancias
desde \texttt{renderer.ts} (n\'umeros 1--3 sobre las flechas superiores) y su
efecto sobre los subdirectorios de \texttt{data/instances/}. Las flechas
discontinuas inferiores indican el vaciado realizado por \texttt{clean.py}
cuando el interruptor de limpieza est\'a activo.}
\label{fig:cap3-flujo}
\end{figure}

Cada invocaci\'on a un script se realiza a trav\'es de la funci\'on auxiliar
\texttt{runPythonCollect}, que envuelve \texttt{window.api.runPython} en
una promesa: se suscribe a los eventos \texttt{python-stdout},
\texttt{python-stderr} y \texttt{python-exit} descritos en el
Cap\'itulo~\ref{cap:1}, acumula la salida completa del proceso y resuelve
la promesa \'unicamente cuando el subproceso finaliza, entregando el
c\'odigo de salida junto con la salida est\'andar y de error completas. Este
patr\'on de sincronizaci\'on es lo que permite que \texttt{downloadSources}
exprese la secuencia de m\'ultiples scripts como una funci\'on
\texttt{async} lineal, mediante sucesivas expresiones \texttt{await}, en
lugar de anidar callbacks.
